\chapter{Relações constitutivas}\label{ch:Constitutiva}
\section{Material elástico}
Considere o modelo de material elástico não linear de St. Venant–Kirchhoff, onde a densidade de energia de deformação $W(E)$ assume a seguinte forma:

\begin{equation}\label{eq:DensidadeEnergia}
	W(\mathbf{E}) = \frac{1}{2} \mathbf{E} : \mathbf{D} : \mathbf{E}.
\end{equation}

A notação $:$ representa o operador de contração para tensores, onde$\mathbf{a}:\mathbf{b} = a_{ij} b_{ij}$ (com soma sobre índices repetidos), e $\mathbf{D}$ é o tensor constitutivo de quarta ordem. Para materiais isotrópicos, definido como:
\begin{equation}
	\mathbf{D} =\lambda \mathbf{1} \otimes \mathbf{1} + 2 \mu \mathbf{I}
\end{equation}

\begin{equation}\label{eq:RelConst}
	D_{ijkl} = \lambda \delta_{ij} \delta_{kl} + \mu (\delta_{ik} \delta_{jl} + \delta_{il} \delta_{jk})
\end{equation}

Formulação utilizada para materiais elásticos lineares. Na Eq.~\eqref{eq:RelConst},$\lambda$ e $\mu$ são as constantes de Lamé para materiais isotrópicos, onde $\delta$ é o delta de Kronecker.$\mathbf{1}$ representa o tensor unitário de segunda ordem, e $I$ é o tensor unitário simétrico de quarta ordem, definido como:

\begin{equation}
	I_{ijkl} = \frac{1}{2} (\delta_{ik} \delta_{jl} + \delta_{il} \delta_{jk}),
\end{equation}

as constantes de Lamé $\lambda$ e $\mu$ para materiais isotrópicos também podem ser expressas em termos do módulo de elasticidade $E$ e do coeficiente de Poisson $\nu$ como:

\begin{equation}
	\lambda = \frac{\nu E}{(1 + \nu)(1 - 2\nu)}, \quad \mu = \frac{E}{2(1 + \nu)}
\end{equation}

A relação constitutiva pode ser obtida diferenciando a Eq.~\eqref{eq:DensidadeEnergia} em relação à deformação Lagrangiana $\mathbf{E}$, resultando em:

\begin{equation}\label{eq:tensaodeformacaolinear}
	\mathbf{S} = \frac{\partial W(\mathbf{E})}{\partial \mathbf{E}} = \mathbf{D} : \mathbf{E} = \lambda \text{tr}(\mathbf{E})\mathbf{1} + 2\mu \mathbf{E}
\end{equation}
onde $S$ representa a tensão de Piola-Kirchhoff de segunda ordem. Note que a relação entre tensão $S$ e deformação $E$ é linear. 

No entanto, muitos materiais desviam-se deste comportamento tensão-deformação linear, especialmente em grandes deformações. Consequentemente, a Eq.~\eqref{eq:tensaodeformacaolinear} é mais adequada para pequenas deformações.
\section{Material Plástico}